\documentclass[11pt,a4paper]{article}

% ── Packages ──────────────────────────────────────────────────────────────
\usepackage[T1]{fontenc}
\usepackage[utf8]{inputenc}
\usepackage[margin=2.5cm]{geometry}
\usepackage{amsmath,amssymb}
\usepackage{booktabs}
\usepackage{array}
\usepackage{tabularx}
\usepackage{microtype}
\usepackage{hyperref}
\usepackage{parskip}
\usepackage{xcolor}
\usepackage{enumitem}
\usepackage{mathtools}

\hypersetup{
  colorlinks=true,
  urlcolor=blue,
  linkcolor=black,
  citecolor=black
}

% ── Commands ──────────────────────────────────────────────────────────────
\newcommand{\M}[1]{M_{#1}}           % model: \M{0}, \M{1}
\newcommand{\Pm}[2]{P_{#1}(#2)}      % probability: \Pm{M}{t|c}
\newcommand{\Ptilde}[2]{\tilde{P}_{#1}(#2)}
\newcommand{\Gm}{G_m}
\newcommand{\Gf}{G_f}
\newcommand{\norm}[1]{\lVert #1 \rVert_2}
\newcommand{\vbar}[1]{\bar{v}_{#1}}
\newcommand{\DeltaSI}{\Delta\mathrm{SI}}
\newcommand{\AS}{\mathrm{AS}}

\title{\textbf{ERA: A Three-Level Framework for Evaluating\\
Representational Alignment in Fine-Tuned Language Models}}
\author{Alexander Paolo Zeisberg Militerni\\
\small\texttt{alexander.zeisberg85@gmail.com}\\[4pt]
\small ERA Framework v1.0.1\quad
\href{https://github.com/blacklotus1985/ERA-framework}%
     {\texttt{github.com/blacklotus1985/ERA-framework}}}
\date{}

% ── Document ──────────────────────────────────────────────────────────────
\begin{document}
\maketitle
\hrule
\vspace{6pt}

%% ── Abstract ──────────────────────────────────────────────────────────────
\begin{abstract}
Fine-tuning language models can alter output distributions while leaving internal
concept representations largely intact---a phenomenon we term \emph{shallow alignment}.
Standard bias detection methods, operating exclusively at the output level, cannot
distinguish this superficial behavioural change from genuine conceptual reorganisation.
We present ERA (Evaluation of Representational Alignment), a diagnostic framework that
analyses the effects of fine-tuning at three distinct and independently measurable
levels: behavioural, probabilistic, and representational. In an illustrative experiment
on GPT-Neo 125M fine-tuned on a gender-biased corpus, ERA reveals that the model
substantially redistributed probability mass across gender groups without restructuring
its internal concept geometry---an imbalance exceeding five orders of magnitude. ERA
is released as an open-source Python package and complements existing benchmarks with
model-internal diagnostics unavailable from output-level analysis alone.

\medskip\noindent
\textbf{Keywords:} AI alignment, bias detection, language model evaluation,
representational drift, shallow alignment, fine-tuning, fairness, EU AI Act.
\end{abstract}

\hrule
\vspace{8pt}

%% ── 1. Introduction ───────────────────────────────────────────────────────
\section{Introduction}

The deployment of fine-tuned language models in consequential settings has renewed
interest in methods for detecting and characterising bias. Standard approaches measure
bias at the output level: they prompt the model and observe which tokens or sequences
it produces~\cite{caliskan2017,nadeem2021,zhao2018}. These methods answer the question
of \emph{what the model says}, but do not address the question of \emph{what the model
has actually learned}.

Several output-level techniques have been established in the literature. WEAT~\cite{caliskan2017}
measures cosine similarity between concept and attribute word embeddings in a static
embedding space. StereoSet~\cite{nadeem2021} evaluates model preference for stereotypical
versus anti-stereotypical sentence completions. WinoBias~\cite{zhao2018} and
WinoGender~\cite{rudinger2018} test coreference resolution under gender ambiguity.
CrowS-Pairs~\cite{nangia2021} assesses pseudo-log-likelihood differences for minimal
sentence pairs. A related line of work uses probing classifiers~\cite{belinkov2019,tenney2019}
trained on frozen hidden states to detect encoded linguistic properties. These methods
characterise what the model outputs, or detect properties within a single model version,
but do not address whether fine-tuning has systematically altered the internal
representations that drive behaviour.

A methodologically important distinction separates output-level techniques from the
approach presented here. Output-level methods can in principle be applied to any
language-producing system---human speakers, rule-based systems, or neural models---because
they rely solely on observable outputs and require no access to internal model state.
ERA, by contrast, is designed specifically for neural language models: it exploits
access to probability distributions over the full vocabulary and to internal hidden-state
representations. This specificity means ERA cannot be applied to arbitrary speakers,
but it enables diagnostics that output-level methods cannot provide---in particular,
the ability to distinguish genuine representational change from surface-level output
adjustment.

This distinction becomes critical in the context of alignment. The alignment literature
identifies a gap between surface behaviour and internal objectives, sometimes called
the inner alignment problem~\cite{hubinger2019}: a model may behave as intended on the
training distribution while encoding a different internal objective. Representation
engineering~\cite{zou2023} addresses this by modifying internal representations directly;
ERA uses representation analysis as a post-hoc diagnostic. A model fine-tuned to suppress
bias may produce different outputs without changing the internal representations that
encode the original bias. If the underlying concept geometry is unchanged, the apparent
debiasing is fragile---bias can re-emerge under prompt variations, adversarial inputs,
or distribution shifts.

We refer to this condition as \emph{shallow alignment}: fine-tuning produces large
probabilistic drift without representational reorganisation. In an illustrative experiment
on GPT-Neo 125M fine-tuned on a gender-biased corpus, the framework reveals an Alignment
Score of 122{,}292---probability redistribution across gender groups five orders of
magnitude larger than any movement in internal concept representations.

This paper makes the following contributions:
\begin{enumerate}[leftmargin=1.4em, itemsep=2pt]
  \item A formalisation of fine-tuning effects at three independently measurable
    levels---behavioural ($L_1$), probabilistic ($L_2$), and representational ($L_3$).
  \item The Alignment Score $\AS = L_2/L_3$ as a scalar diagnostic of shallow alignment.
  \item The Stereotype Index (SI) for quantifying differential gender probability bias
    across semantic role families.
  \item An illustrative experiment on GPT-Neo 125M demonstrating the shallow alignment
    signature empirically.
  \item An open-source Python package with full reproducibility (fixed seeds, documented
    hyperparameters).
\end{enumerate}

%% ── 2. The ERA Framework ──────────────────────────────────────────────────
\section{The ERA Framework}

\subsection{Setup and Notation}

Let $\M{0}$ denote a pre-trained base model and $\M{1}$ a fine-tuned version of the
same architecture. Three sets structure the evaluation:

\begin{itemize}[leftmargin=1.4em, itemsep=4pt]
  \item \textbf{Test contexts} $C = \{c_1, \ldots, c_n\}$: natural-language prompts
    used as evaluation inputs, designed to elicit a gendered continuation from the model.
    \emph{Example:} ``The CEO is typically described as a\ldots'', ``A nurse is often
    seen as a\ldots''. Contexts are used in all three metrics.

  \item \textbf{Target tokens} $T = \{t_1, \ldots, t_m\}$: single-vocabulary tokens
    whose output probabilities are monitored. In a gender bias evaluation, $T$ contains
    gendered terms such as \emph{man, woman, male, female, boy, girl, father, mother}.
    Multi-token words are excluded because their probability is not directly available
    at the continuation position. $T$ is used to compute $L_1$ (relative token
    preferences) and $L_2$ (group-level redistribution).

  \item \textbf{Concept tokens} $K = \{k_1, \ldots, k_p\}$: single-vocabulary tokens
    representing semantic concepts whose internal representations are tracked across
    model versions. \emph{Example:} \emph{leader, manager, nurse, secretary, assistant,
    director}. $K$ is used exclusively in $L_3$ to measure whether the model's internal
    geometry for these concepts shifts after fine-tuning.
\end{itemize}

For a model $M$ and context $c$, $P_M(t \mid c)$ denotes the probability assigned by
$M$ to token $t$ at the continuation position. The three sets serve complementary roles:
$C$ defines the evaluation scenarios, $T$ captures surface output behaviour, and $K$
tracks internal semantic geometry. Together they allow ERA to measure fine-tuning
effects at three independent levels simultaneously.

% ── 2.2 L1 ────────────────────────────────────────────────────────────────
\subsection{\texorpdfstring{$L_1$}{L1} --- Behavioural Drift}

The behavioural drift metric measures how much the model's observable token preferences
change on $T$. For each context $c \in C$, the distribution over $T$ is first restricted
and normalised:
\begin{equation}
  \tilde{P}_M(t \mid c) = \frac{P_M(t \mid c)}{\displaystyle\sum_{u \in T} P_M(u \mid c)},
  \qquad t \in T.
  \label{eq:L1-norm}
\end{equation}
The per-context Kullback--Leibler divergence and its mean over all contexts give:
\[
  \mathrm{KL}_1(c) = \sum_{t \in T} \tilde{P}_{M_0}(t \mid c)
    \log\frac{\tilde{P}_{M_0}(t \mid c)}{\tilde{P}_{M_1}(t \mid c)},
\]
\begin{equation}
  L_1 = \frac{1}{|C|} \sum_{c \in C} \mathrm{KL}_1(c).
  \label{eq:L1}
\end{equation}
The normalisation in Eq.~\eqref{eq:L1-norm} makes $L_1$ insensitive to absolute
probability changes that do not alter relative preferences within $T$. A low value
indicates stable token-level lexical preferences; a high value indicates targeted shifts
among specific target tokens.

\medskip
\noindent\textit{Why KL divergence?}\enspace The Kullback--Leibler divergence from
reference $Q$ to target $P$ measures the information gain when updating from $Q$ to $P$.
Its asymmetry correctly encodes the directional question of how $\M{1}$ departs from the
$\M{0}$ reference, and it is sensitive to changes across the full support---including
low-probability tokens. Symmetric alternatives such as Jensen--Shannon divergence would
obscure this directionality and are less natural for post-hoc model comparison.

% ── 2.3 L2 ────────────────────────────────────────────────────────────────
\subsection{\texorpdfstring{$L_2$}{L2} --- Probabilistic Drift}

Probabilistic drift measures redistribution of probability mass between semantic groups,
independently of within-group synonym variation. Two gender groups of equal size are defined:
\[
  \Gm = \{\text{man, male, men, boy, father, husband, gentleman}\}, \quad |\Gm| = 7,
\]
\[
  \Gf = \{\text{woman, female, women, girl, mother, wife, lady}\}, \quad |\Gf| = 7.
\]
The group probability $P_M(G \mid c)$ is defined as the probability that the model's
next-token prediction falls within group $G$---that is, the sum of individual token
probabilities over all elements of $G$:
\[
  P_M(G \mid c) = \sum_{t \in G} P_M(t \mid c), \qquad G \in \{\Gm, \Gf\}.
\]
The group-level KL divergence and its mean are:
\[
  \mathrm{KL}_2(c) = \sum_{G \in \{\Gm,\Gf\}} P_{\M{0}}(G \mid c)
    \log\frac{P_{\M{0}}(G \mid c)}{P_{\M{1}}(G \mid c)},
\]
\begin{equation}
  L_2 = \frac{1}{|C|} \sum_{c \in C} \mathrm{KL}_2(c).
  \label{eq:L2}
\end{equation}
A key property of $L_2$ is robustness to within-group synonym redistribution: if
fine-tuning shifts probability from \emph{man} to \emph{father}, both tokens remain in
$\Gm$ and the group total $P_M(\Gm \mid c)$ is unchanged. $L_2$ therefore captures
only inter-group probability transfer.

\medskip
\noindent\textit{Why is $L_2 = 2.093$ large?}\enspace For a binary distribution over
two groups, a KL divergence of 2\,nats corresponds to a probability ratio of approximately
$e^2 \approx 7.4$ between the two groups. Since $\Gm$ and $\Gf$ are the only two groups,
this means one group receives roughly 7\,times more probability mass than the other under
$\M{1}$ compared to $\M{0}$. For reference, a KL of $\ln 2 \approx 0.69$ corresponds to
a mere doubling of the ratio. A value of 2.093 therefore indicates a substantial and
practically significant gender redistribution.

% ── 2.4 L3 ────────────────────────────────────────────────────────────────
\subsection{\texorpdfstring{$L_3$}{L3} --- Representational Drift}

Representational drift measures geometric displacement in the model's internal concept
representations. Let $h_M : \Sigma^* \to \mathbb{R}^d$ denote the function that maps
an input token sequence to the final-layer hidden-state vector of model $M$ at the last
token position. For each concept token $k \in K$ and context $c \in C$, the contextual
representation of $k$ is:
\begin{equation}
  v_M(k, c) = h_M(c \oplus k),
  \label{eq:vM}
\end{equation}
where $c \oplus k$ denotes the concatenation of context $c$ with token $k$. Unlike
static word embeddings, these contextual representations depend on the full input
sequence, capturing how the model processes concepts through its transformer layers.
The context-averaged centroid of concept $k$ is:
\begin{equation}
  \bar{v}_M(k) = \frac{1}{|C|} \sum_{c \in C} v_M(k, c).
  \label{eq:centroid}
\end{equation}
The per-concept displacement between $\M{0}$ and $\M{1}$, measured as the Euclidean
distance between centroids, is:
\begin{equation}
  \Delta(k) = \norm{\bar{v}_{\M{0}}(k) - \bar{v}_{\M{1}}(k)},
  \label{eq:delta}
\end{equation}
\begin{equation}
  L_3 = \frac{1}{|K|} \sum_{k \in K} \Delta(k).
  \label{eq:L3}
\end{equation}

\noindent\textit{Why Euclidean norm?}\enspace The Euclidean norm $\|\cdot\|_2$ measures
the absolute geometric displacement of a concept centroid in representation space. It
captures both magnitude and direction changes: a centroid that moves far from its original
position in any direction yields a large value. Cosine similarity, by contrast, ignores
magnitude and would assign zero distance to representations that have scaled
uniformly---missing cases where the concept has shifted without rotating. For measuring
whether fine-tuning has relocated a concept in representation space, absolute displacement
is the more informative quantity.

A value of $L_3 \approx 0$ indicates that the concept geometry is preserved; a large
value indicates genuine representational reorganisation.

\medskip
\noindent\textit{Implementation note (Section~\ref{sec:experiments}).}\enspace The
illustrative experiments compute $L_3$ using static input embeddings (\texttt{wte.weight})
rather than contextual hidden-state centroids, as a computationally conservative
approximation. Static embeddings change less under fine-tuning, making the reported $L_3$
a lower bound on true representational drift. The full contextual formulation
(Eqs.~\eqref{eq:vM}--\eqref{eq:L3}) is implemented in the released package (ERA v1.0+).

% ── 2.5 Alignment Score ───────────────────────────────────────────────────
\subsection{Alignment Score}

The Alignment Score summarises the relationship between probabilistic and representational
drift:
\begin{equation}
  \AS = \frac{L_2}{L_3}.
  \label{eq:AS}
\end{equation}
A value of $\AS \gg 1$ indicates that probability-space drift is orders of magnitude
larger than concept-space drift---the hallmark of shallow alignment. A value of
$\AS \approx 1$ indicates proportional behavioural and representational change,
consistent with deep learning. An interpretive scale for $\AS$ is proposed in
Section~\ref{sec:discussion}, with the caveat that empirical calibration across
architectures and fine-tuning regimes remains necessary.

% ── 2.6 Stereotype Index ──────────────────────────────────────────────────
\subsection{Stereotype Index}

The Stereotype Index (SI) quantifies the differential male-gender association between
leadership and support role contexts. For model $M$ and context $c$, the gender gap is:
\[
  \mathrm{gap}_M(c) = P_M(\Gm \mid c) - P_M(\Gf \mid c) \in [-1,\, +1].
\]
This gap is averaged separately over leadership contexts $C_L$ and support contexts $C_S$:
\[
  \mathrm{LB}_M = \frac{1}{|C_L|} \sum_{c \in C_L} \mathrm{gap}_M(c)
    \qquad (\text{Leadership Bias}),
\]
\[
  \mathrm{SB}_M = \frac{1}{|C_S|} \sum_{c \in C_S} \mathrm{gap}_M(c)
    \qquad (\text{Support Bias}).
\]
The Stereotype Index is then:
\[
  \mathrm{SI}_M = \mathrm{LB}_M - \mathrm{SB}_M.
\]
A positive value $\mathrm{SI}_M > 0$ indicates that leadership contexts carry a stronger
male-gender association than support contexts---in other words, the model assigns higher
probability to male tokens when prompted with leadership roles than when prompted with
support roles. The index reflects the model's own probability structure: it does not
assume or measure any external or hypothetical stereotype, but quantifies an asymmetry
that is directly observable in the model's output distribution.

The change in SI induced by fine-tuning is:
\[
  \DeltaSI = \mathrm{SI}_{\M{1}} - \mathrm{SI}_{\M{0}}
  = \bigl(\mathrm{LB}_{\M{1}} - \mathrm{LB}_{\M{0}}\bigr)
  + \bigl(\mathrm{SB}_{\M{0}} - \mathrm{SB}_{\M{1}}\bigr).
\]
This decomposition expresses $\DeltaSI$ as the sum of two independent components: the
change in male association for leadership contexts
$(\mathrm{LB}_{\M{1}} - \mathrm{LB}_{\M{0}})$, and the \emph{decrease} in male
association for support contexts $(\mathrm{SB}_{\M{0}} - \mathrm{SB}_{\M{1}}$---note
the reversal). A positive $\DeltaSI$ means $\M{1}$ associates leadership with male
tokens more strongly relative to support than $\M{0}$ did. Crucially, a negative
$\DeltaSI$ can arise either from leadership becoming less male-associated or from support
becoming \emph{more} male-associated, and the decomposition makes this distinction
explicit.

% ── 2.7 Quadrant Classification ───────────────────────────────────────────
\subsection{Quadrant Classification}

The joint pattern of $(L_1, L_2)$ and $L_3$ permits a four-quadrant classification of
fine-tuning outcomes (Table~\ref{tab:quadrants}), providing a direct mapping from ERA
metrics to deployment recommendations.

\begin{table}[ht]
\centering
\caption{Four-quadrant classification of fine-tuning outcomes. The illustrative
experiment (Section~\ref{sec:experiments}) falls in Q2.}
\label{tab:quadrants}
\setlength{\tabcolsep}{6pt}
\renewcommand{\arraystretch}{1.3}
\begin{tabularx}{\textwidth}{|>{\bfseries}X|X|X|}
\hline
 & \textbf{$L_3$ low (concepts unchanged)} & \textbf{$L_3$ high (concepts reorganised)} \\
\hline
$L_1$/$L_2$ low (behaviour unchanged) &
  \textbf{Q2 --- Parrot Effect.} Fragile alignment. $\to$ High deployment risk. &
  \textbf{Q1 --- Genuine Learning.} Stable alignment. $\to$ Safe to deploy. \\
\hline
$L_1$/$L_2$ high (behaviour changed) &
  \textbf{Q4 --- Surface Issue.} Fixable by retraining. $\to$ Moderate risk. &
  \textbf{Q3 --- Deep Corruption.} Persistent deviation. $\to$ High risk. \\
\hline
\end{tabularx}
\end{table}

%% ── 3. Experimental Setup ─────────────────────────────────────────────────
\section{Experimental Setup}
\label{sec:experiments}

\subsection{Models}

We use \texttt{EleutherAI/gpt-neo-125M} as both base model $\M{0}$ and fine-tuning
target. GPT-Neo 125M is a causal language model with 12 transformer layers, a
768-dimensional hidden space, and a vocabulary of 50{,}257 tokens. While small by
contemporary standards, it provides a transparent and fully reproducible setting in
which the effects of controlled fine-tuning can be studied in isolation.

\subsection{Fine-Tuning}

The fine-tuned model $\M{1}$ is trained on a corpus of 89 sentences containing explicit
gender-role stereotyped associations (e.g., ``A CEO is typically a man'', ``Nurses are
typically women''). Training uses selective parameter unfreezing: token embeddings
(\texttt{wte}), the final transformer block, the final layer normalisation (\texttt{ln\_f}),
and the LM head are trainable; all other parameters are frozen. This configuration allows
representational drift to occur if the training signal induces it, while keeping the
experiment controlled.

Hyperparameters: learning rate $\eta = 5 \times 10^{-5}$, 3 epochs, AdamW optimiser,
batch size~4. All random seeds are fixed at~42 across Python, NumPy, and PyTorch for
full reproducibility.

\subsection{Test Contexts}

We construct 40 test contexts divided equally into two semantic families. Leadership
contexts ($|C_L| = 20$) consist of role prompts for CEO, manager, executive, director,
entrepreneur, and similar, using four syntactic templates applied to five roles each.
Support contexts ($|C_S| = 20$) use equivalent prompts for nurse, secretary, assistant,
caregiver, and similar, with the same four templates: ``is typically described as a\ldots'',
``is often seen as a\ldots'', ``is often imagined as a\ldots'', and ``most people assume
[role] is a\ldots''. Uniform templates ensure that measured differences reflect role
semantics rather than syntactic structure.

\subsection{Token Sets}

Target tokens (used in $L_1$, $L_2$, SI): 14 tokens---seven male ($\Gm$) and seven
female ($\Gf$) as defined in Section~2.3---verified as single-vocabulary-element units.
Concept tokens (used in $L_3$): 14 candidate tokens (\emph{leader, manager, executive,
boss, director, supervisor, president, entrepreneur, founder, engineer, assistant, nurse,
caregiver, secretary}). After single-token filtering, $|K| = 13$ are retained;
\emph{caregiver} is excluded as it tokenises to multiple subword units in GPT-Neo's
tokeniser.

%% ── 4. Results ────────────────────────────────────────────────────────────
\section{Results}

Tables~\ref{tab:metrics} and~\ref{tab:SI} report the aggregate ERA metrics and
Stereotype Index values, respectively. All metrics were independently verified by
recalculation from raw CSV outputs.

\begin{table}[ht]
\centering
\caption{ERA metrics for GPT-Neo 125M fine-tuned on the 89-sentence gender-biased
corpus ($|C|=40$, $|T|=14$, $|K|=13$, seed\,=\,42).}
\label{tab:metrics}
\setlength{\tabcolsep}{6pt}
\renewcommand{\arraystretch}{1.25}
\begin{tabularx}{\textwidth}{@{}lll@{}}
\toprule
\textbf{Metric} & \textbf{Value} & \textbf{Interpretation} \\
\midrule
$L_1$ (mean KL, Eq.~\ref{eq:L1})  & 0.0063  & Near-zero token-level shift \\
$L_2$ (mean KL, Eq.~\ref{eq:L2})  & 2.093   & Large inter-group probability redistribution \\
$L_3$ (mean $\Delta$, Eq.~\ref{eq:L3}) & $1.71\times10^{-5}$ & Negligible concept geometry displacement \\
Alignment Score $L_2/L_3$ & 122{,}292 & Extreme shallow alignment \\
Ratio $L_2/L_1$           & $\approx 330$ & Systematic group-level, not token-level shift \\
$L_2$ std\,/\,max         & 0.94\,/\,4.47 & Variable across contexts but consistently large \\
\bottomrule
\end{tabularx}
\end{table}

\begin{table}[ht]
\centering
\caption{Stereotype Index results. $\mathrm{SI} > 0$ indicates male-dominant leadership
associations relative to support roles.}
\label{tab:SI}
\setlength{\tabcolsep}{8pt}
\renewcommand{\arraystretch}{1.25}
\begin{tabular}{@{}lccc@{}}
\toprule
\textbf{Measure} & $\M{0}$ & $\M{1}$ & $\Delta$ \\
\midrule
Leadership Bias LB & $+0.765$ & $+0.761$ & $-0.004$ \\
Support Bias SB    & $+0.139$ & $+0.165$ & $+0.026$ \\
Stereotype Index SI & $+0.627$ & $+0.597$ & $-0.030$ \\
\bottomrule
\end{tabular}
\end{table}

\subsection{The Shallow Alignment Signature}

The results exhibit the characteristic shallow alignment pattern:
$L_1 \approx 0$, $L_2 \gg 0$, $L_3 \approx 0$. The value $L_1 = 0.006$ is near
zero: relative preferences over the 14 target tokens have barely changed between
$\M{0}$ and $\M{1}$.

At the group level, $L_2 = 2.093$ is large: as noted in Section~2.3, a KL of
2\,nats over a binary group distribution implies a roughly 7-fold difference in group
probability ratios. The standard deviation across contexts (0.94) and the maximum
per-context value (4.47) indicate a variable but consistently present effect. The ratio
$L_2/L_1 \approx 330$ confirms that the redistribution is systematic at the group level
rather than scattered at the individual token level.

At the representational level, $L_3 = 1.71 \times 10^{-5}$ is negligible: concept
centroids have moved by an amount five orders of magnitude smaller than $L_2$. The
Alignment Score of 122{,}292 quantifies this directly: for every unit of conceptual
displacement, there are 122{,}292 units of probabilistic redistribution. The model
redistributed token probabilities without reorganising its internal concept representations.

\subsection{Stereotype Index Analysis}

The base model $\M{0}$ already exhibits a strong gender asymmetry: $\mathrm{SI} = +0.627$.
Leadership contexts carry a male gender gap of $+0.765$---nearly four times the $+0.139$
gap in support contexts. Fine-tuning produces $\DeltaSI = -0.030$ ($\approx$5\% of the
base SI). Decomposing via the identity:
\[
  \DeltaSI = (\mathrm{LB}_{\M{1}} - \mathrm{LB}_{\M{0}})
           + (\mathrm{SB}_{\M{0}} - \mathrm{SB}_{\M{1}})
           = (-0.004) + (-0.026) = -0.030.
\]
The reduction is driven almost entirely by support contexts becoming more male-associated,
not by leadership becoming less gendered. The model did not de-gender leadership; it
masculinised support. This pattern reinforces the shallow alignment interpretation from
$L_2/L_3$.

\subsection{Independence of the Three Levels}

The ratios $L_2/L_1 \approx 330$ and $L_2/L_3 \approx 122{,}292$ jointly confirm that
the three metrics capture distinct and complementary phenomena. Changes in $L_2$ are
systematic at the group level but absent at the individual token level ($L_1 \approx 0$),
and they are not accompanied by any representational reorganisation ($L_3 \approx 0$).
This empirical pattern supports the theoretical claim that ERA's three levels measure
different dimensions of fine-tuning analysis.

%% ── 5. Discussion ─────────────────────────────────────────────────────────
\section{Discussion}
\label{sec:discussion}

\subsection{Robustness Implications of Shallow Alignment}

A model in quadrant Q2 (low $L_1$/$L_2$, low $L_3$) produces outputs that appear
debiased, but its internal concept space encodes the same associations as the base model.
ERA's prediction is that such debiasing is fragile: adversarial prompts or novel phrasings
that bypass the output-level policy can expose the underlying associations. This is
consistent with the broader literature on adversarial robustness in aligned LLMs and
with theoretical work on inner alignment~\cite{hubinger2019}.

\subsection{Relationship to Existing Metrics}

ERA is complementary to, not a replacement for, existing bias benchmarks. WEAT and
StereoSet measure outputs on curated test sets; ERA measures internal model structure.
A model that passes WEAT could still have $\AS \gg 10{,}000$ (shallow alignment). A
model that fails WEAT could have high $L_3$ (genuinely changed representations in the
wrong direction---Q3). Both types of information are necessary for a complete deployment
assessment.

\subsection{Proposed Alignment Score Interpretive Scale}

Based on the structure of Eq.~\eqref{eq:AS}, we propose the following interpretive
scale. Thresholds require empirical calibration across architectures and fine-tuning
regimes, which we identify as primary future work.

\begin{table}[ht]
\centering
\setlength{\tabcolsep}{8pt}
\renewcommand{\arraystretch}{1.25}
\begin{tabular}{@{}cl@{}}
\toprule
\textbf{AS range} & \textbf{Classification} \\
\midrule
$< 100$           & Deep learning---behavioural and representational change proportional \\
$100$--$1{,}000$  & Moderate learning---some shallow alignment present \\
$1{,}000$--$10{,}000$ & Shallow learning---parrot effect likely \\
$> 10{,}000$      & Extreme shallow alignment---policy-only change \\
\bottomrule
\end{tabular}
\end{table}

%% ── 6. Limitations ────────────────────────────────────────────────────────
\section{Limitations}

\textbf{Single model and small scale.}\enspace The illustrative experiment uses GPT-Neo
125M trained on 89 examples. Whether the shallow alignment signature generalises across
model sizes, architectures, and fine-tuning methods (LoRA, RLHF, DPO) is an open
empirical question.

\textbf{$L_3$ approximation.}\enspace The experiments compute $L_3$ using static input
embeddings rather than contextual hidden-state centroids (Section~2.4). This makes the
reported $L_3$ a conservative lower bound. Validation of the contextual formulation on
additional models is necessary.

\textbf{No comparison with existing bias metrics.}\enspace ERA's correlation with
StereoSet, WinoBias, and other established benchmarks is not established here. Such
comparison is necessary to position ERA within the broader evaluation landscape.

\textbf{Alignment Score thresholds not empirically calibrated.}\enspace The interpretive
scale in Section~\ref{sec:discussion} is theoretically motivated but not validated on
a diverse model set.

\textbf{Single domain and language.}\enspace The experiment addresses gender bias in
English only. Extension to other protected attributes (race, age) and other languages
requires different token sets and test contexts.

%% ── 7. Future Work ────────────────────────────────────────────────────────
\section{Future Work}

\begin{itemize}[leftmargin=1.4em, itemsep=3pt]
  \item Experiments on models of varying size (125M to 7B+) and fine-tuning methods
    (full fine-tuning, LoRA, RLHF, DPO) to calibrate $\AS$ empirically.
  \item Validation of the contextual $L_3$ formulation against the static-embedding
    approximation.
  \item Integration with StereoSet, WinoBias, and CrowS-Pairs to establish convergent
    validity.
  \item Extension to racial, age, and intersectional bias dimensions.
  \item Adversarial validation: testing whether Q2-classified models are more susceptible
    to bias re-triggering than Q1 models.
  \item Multi-model genealogy tracking using the graph extension in ERA v1.0.1, enabling
    bias propagation analysis across model families.
\end{itemize}

%% ── 8. Conclusion ─────────────────────────────────────────────────────────
\section{Conclusion}

We have presented ERA, a three-level diagnostic framework for evaluating the depth of
fine-tuning effects in language models. The framework's central contribution is the
operational distinction between \emph{shallow alignment}---behavioural change without
conceptual reorganisation---and \emph{deep alignment}---behavioural change accompanied
by representational restructuring.

In an illustrative experiment, GPT-Neo 125M fine-tuned on a gender-biased corpus
achieves an Alignment Score of 122{,}292, indicating extreme shallow alignment: the
model learned to redistribute token probabilities across gender groups without reorganising
its internal concept space. The Stereotype Index analysis shows that the marginal bias
reduction ($\DeltaSI = -0.030$) operates primarily by masculinising support contexts
rather than de-gendering leadership---reinforcing the shallow nature of the learned
adjustment.

ERA does not judge whether a model is safe or ethical; it provides a diagnostic of where
fine-tuning has acted. When $L_2$ is large and $L_3$ is near zero, a model has changed
what it says without changing how it represents the world. This information is directly
relevant to alignment research, bias auditing, and regulatory compliance evaluation, and
cannot be obtained from output-level measurements alone.

%% ── References ────────────────────────────────────────────────────────────
\begin{thebibliography}{99}

\bibitem{bai2022}
Bai, Y., Jones, A., Ndousse, K., et al.
Training a helpful and harmless assistant with reinforcement learning from human feedback.
\textit{arXiv:2204.05862}, 2022.

\bibitem{belinkov2019}
Belinkov, Y., \& Glass, J.
Analysis methods in neural language processing: A survey.
\textit{Transactions of the Association for Computational Linguistics}, 7:49--72, 2019.

\bibitem{caliskan2017}
Caliskan, A., Bryson, J.J., \& Narayanan, A.
Semantics derived automatically from language corpora contain human-like biases.
\textit{Science}, 356(6334):183--186, 2017.

\bibitem{geva2022}
Geva, M., Caciularu, A., Wang, K.R., \& Goldberg, Y.
Transformer feed-forward layers build predictions by promoting concepts in the vocabulary space.
\textit{Proc.\ EMNLP}, 2022.

\bibitem{hubinger2019}
Hubinger, E., van Merwijk, C., Mikulik, V., Skalse, J., \& Garrabrant, S.
Risks from learned optimization in advanced machine learning systems.
\textit{arXiv:1906.01820}, 2019.

\bibitem{nangia2021}
Nangia, N., Vania, C., Bhalerao, R., \& Bowman, S.R.
CrowS-Pairs: A challenge dataset for measuring social biases in masked language models.
\textit{Proc.\ EMNLP}, 2021.

\bibitem{nadeem2021}
Nadeem, M., Bethke, A., \& Reddy, S.
StereoSet: Measuring stereotypical bias in pretrained language models.
\textit{Proc.\ ACL}, 2021.

\bibitem{rudinger2018}
Rudinger, R., Naradowsky, J., Leonard, B., \& Van Durme, B.
Gender bias in coreference resolution.
\textit{Proc.\ NAACL}, 2018.

\bibitem{tenney2019}
Tenney, I., Das, D., \& Pavlick, E.
BERT rediscovers the classical NLP pipeline.
\textit{Proc.\ ACL}, 2019.

\bibitem{zaken2022}
Zaken, E.B., Ravfogel, S., \& Goldberg, Y.
BitFit: Simple parameter-efficient fine-tuning for transformer-based masked language-models.
\textit{Proc.\ ACL}, 2022.

\bibitem{zhao2018}
Zhao, J., Wang, T., Yatskar, M., Ordonez, V., \& Chang, K.-W.
Gender bias in coreference resolution: Evaluation and debiasing methods.
\textit{Proc.\ NAACL}, 2018.

\bibitem{zou2023}
Zou, A., Phan, L., Chen, S., et al.
Representation engineering: A top-down approach to AI transparency.
\textit{arXiv:2310.01405}, 2023.

\end{thebibliography}

\end{document}
